\documentclass[11pt]{article}
\usepackage[a4paper,margin=0.82in]{geometry}
\usepackage[T1]{fontenc}
\usepackage[utf8]{inputenc}
\usepackage{lmodern,amsmath,amssymb,amsthm,microtype,booktabs,tabularx,array,xcolor}
\usepackage[numbers]{natbib}
\usepackage[colorlinks=true,linkcolor=blue!55!black,citecolor=blue!55!black,urlcolor=blue!55!black]{hyperref}
\setlength{\emergencystretch}{4em}
\newtheorem{theorem}{Theorem}
\newtheorem{lemma}{Lemma}
\newtheorem{remark}{Remark}
\newcommand{\e}{\mathrm{e}}
\title{\textbf{Rational Atoms after Residue Splitting: Determinant Inflation and the DFT Ledger}}
\author{Liang Wang\\Huazhong University of Science and Technology}
\date{July 2026}
\begin{document}
\maketitle
\begin{abstract}
A rational twist of conductor R splits into R residue-class sums. On each class the two affine slopes have gcd exactly R and the determinant becomes 2R.  One Fourier mode is one DFT coordinate, not control of every residue sum.
The classification is \texttt{ALGEBRAIC\_\allowbreak{}REDUCTION\_\allowbreak{}L1} and the
verdict is \texttt{PROVED\_\allowbreak{}RESIDUE\_\allowbreak{}SPLIT\_\allowbreak{}WITH\_\allowbreak{}DETERMINANT\_\allowbreak{}2R}.  No program-positive L2
claim or endpoint exponent credit is made.
\end{abstract}

\section{Target contract and source boundary}
The target has six simultaneous axes: actual fixed-\(h_0\) packet,
source-locked named physical atom, every deterministic prefix, every
deterministic scale, a fixed-\(X\) power at that atom, and actual active
support.  The value \(h_0=2\) is source-backed data only.  Repository hashes
are integrity locks, not theorem evidence.  TPC-193's declared seven-source
corpus remains \textsc{stop-scoped} \citep{WangTPC193}.


\section{Exact residue decomposition}
Let \(\alpha=r/R\) in lowest terms and let
\[
 C_r(Z)=\sum_{z\in Z}c(z)\e(-rz/R).
\]
Partitioning by \(z=b+Rm\) gives the identity
\begin{equation}\label{eq:dft}
 C_r(Z)=\sum_{b\bmod R}\e(-rb/R)
 \sum_{\substack{m\\b+Rm\in Z}}c(b+Rm).
\end{equation}
For the determinant-two forms \(d+sz,u+az\), the two forms on the
\(b\)-th class are
\[
 (d+sb)+sRm,\qquad (u+ab)+aRm.
\]

\begin{theorem}[Determinant and slope-gcd ledger]
If \(su-ad=2\) and \(as\) is odd, then \((a,s)=1\).  After residue
splitting the two slopes \(sR,aR\) have gcd exactly \(R\), while their
determinant is
\[
 (sR)(u+ab)-(aR)(d+sb)=2R.
\]
Equation \eqref{eq:dft} is one DFT coordinate.  Recovering every residue sum
requires all \(R\) Fourier coordinates and DFT inversion.
\end{theorem}
\begin{proof}
Any common divisor of \(a,s\) divides \(2\); oddness forces it to be one.
The displayed determinant follows by expansion.  The final assertion is
the invertibility of the \(R\times R\) Fourier matrix; a single row has a
kernel of dimension \(R-1\) when \(R>1\).
\end{proof}

\section{Imported identities and new crosswalk}
The DFT identity and the raw progression determinant law are imported from
the TPC-108 alias calculus and TPC-94 phase ledger
\citep{WangTPC94,WangTPC108}.  The L1 contribution here is their joint typed
crosswalk to the TPC-194 decorated direct-prefix object, including the
warning that a single mode is not all residue classes.

\section{Firewall}
The slope-gcd-\(R\), determinant-\(2R\) residue problem is not the original
primitive-slope determinant-two problem.  A theorem for one rational
Fourier mode also cannot be relabelled as simultaneous residue-class
cancellation.


\section{Loss, level and scope ledger}
The smallest missing literal theorem is
\[
\texttt{UNIFORM\_\allowbreak{}ALL\_\allowbreak{}RESIDUE\_\allowbreak{}CLASS\_\allowbreak{}CANCELLATION\_\allowbreak{}OR\_\allowbreak{}DIRECT\_\allowbreak{}MODE\_\allowbreak{}THEOREM}.
\]
The next route is \texttt{RATIONAL\_\allowbreak{}ATOM\_\allowbreak{}THEOREM\_\allowbreak{}SCREEN}.
No new method cell is stopped in this paper.

The named-atom exponent credit is \(0\), while the endpoint requires a
strict budget greater than \(1/400\); the ledger is unpaid.  Block sums are
not cumulative prefixes, symbolic packet atoms are not named production
atoms, phase \(L^2\) and almost-everywhere statements are not pointwise
evaluation, and scoped method failures are not theorem or architecture
failures.

\section{Machine certificate}
The adjacent canonical payload freezes the formula or finite witness,
exact repository-source hashes, any explicitly non-hashed external locator
record, six target axes, scoped stop, route state and claim firewall.  Its
recursive exact schema has closed objects, exact array positions and
constant leaves.  The checker recomputes the finite certificate and executes
ten adversarial mutations.

\bibliographystyle{plainnat}
\bibliography{references}
\end{document}
